\documentclass{report}
\usepackage[utf8]{inputenc}
\usepackage[T1]{fontenc}
\usepackage{geometry}
\usepackage{graphicx}
\usepackage{float}
\usepackage{hyperref}
\usepackage[english,indonesian]{babel}

\title{Optimasi Pemrosesan Template LaTeX Menggunakan Pendekatan Berbasis AI}
\author{Academic Support Agent}
\date{\today}

\begin{document}
\maketitle

\begin{abstract}

\end{abstract}

\chapter{PENDAHULUAN}


\chapter{TINJAUAN PUSTAKA}


% Original content (commented out):
% \documentclass[12pt]{report}
\usepackage[a4paper]{geometry} % Configuración de márgenes y dimensiones de página

\usepackage[utf8]{inputenc} % Codificación de caracteres
\usepackage[spanish]{babel} % Configuración de idioma español
\usepackage{amssymb} % Símbolos matemáticos adicionales
\usepackage{subfigure} % Creación de subfiguras dentro de una figura
\usepackage{newunicodechar} % Soporte para caracteres Unicode
\usepackage{graphicx} % Inclusión de imágenes
\usepackage{amsmath} % Ampliación de funcionalidades matemáticas
\usepackage{xspace} % Control del espacio después de comandos
\usepackage{booktabs} % Mejora la apariencia de las tablas
\usepackage[table,xcdraw]{xcolor} % Colores personalizados para tablas
\usepackage{url} % Inclusión y formato de enlaces URL
\usepackage{lineno} % Numeración de líneas
\usepackage{enumitem} % Control de listas y enumeraciones
\usepackage{soul} % Resaltado de texto con subrayado
\usepackage[numbers]{natbib} % Citas bibliográficas y bibliografía
\usepackage[colorlinks=true, bookmarks=false, citecolor=blue, urlcolor=blue, linkcolor=blue, linktoc=page]{hyperref} % Creación de enlaces y referencias
\usepackage{titlesec} % Control del formato de títulos de sección
\usepackage{fancyhdr} % Personalización de encabezados y pies de página
\usepackage{datetime} % Manipulación de fechas y tiempos
% Modificar márgenes
\geometry{left=2.5cm, right=2.5cm, top=2.0cm, bottom=2.0cm}
% Espaciado entre líneas
\linespread{1.5} 
% Cambiar nombre a tablas
\addto\captionsspanish{\renewcommand{\tablename}{Tabla}}
% Cambiar nombre a lista de tablas
\addto\captionsspanish{\renewcommand{\listtablename}{Índice de tablas}}
% Modificar el estilo de los capítulos
\titleformat{\chapter}[display]
{\color{black}\fontfamily{pcr}\LARGE\scshape\raggedleft}
{\chaptertitlename\ \color{red}\fontsize{20}{00}\mdseries\thechapter} {-10pt} {\color{black}\rmfamily\huge}
%\titlespacing*{\chapter}{0pt}{0pt}{50pt}
% Diseño del encabezado y pie de página
\pagestyle{fancy}
\fancyhf{}
\fancyhead[L]{\nouppercase{\leftmark}}
\fancyhead[R]{\nouppercase{\rightmark}}
\fancyfoot[C]{\thepage}
\renewcommand{\headrulewidth}{0.1pt}
% Quitar espacio vertical en listas
\setlist[itemize]{noitemsep, topsep=0pt}
% Numerar las secciones
\setcounter{secnumdepth}{4} % Establece la profundidad de numeración de secciones
\titleformat{\paragraph}{\normalfont\normalsize\bfseries}{\theparagraph}{1em}{} % Define el formato de las subsubsubsecciones
% Modificar la tabla de contenido
\setcounter{tocdepth}{4}
\addto\captionsspanish{\renewcommand{\contentsname}{Contenido}}


%%%%%%%%%%%%%%%%%%%%%%%%%%% AQUI INICIA EL DOCUMENTO %%%%%%%
\begin{document}
\pagenumbering{Roman}
%%%%%%%%%%%%%%%%%%%%%%%%%%% PORTADA %%%%%%%%%%%%%%%%%%%%%%%%
\begin{titlepage}
\newgeometry{left=2.5cm,right=2.5cm,top=2.5cm,bottom=2.5cm}
\begin{minipage}[c]{0.15\textwidth}
  \centering
  \includegraphics[width=\textwidth]{01_Figuras/00_TECNM.png}
\end{minipage}
\hfill % Espacio horizontal entre las cajas
\begin{minipage}[c]{0.8\textwidth}
    \centering
        \textsc{\large Tecnológico Nacional de México en Celaya}\\[0.3cm]
        \textsc{\large Departamento de Ingeniería Química}\\[0.3cm]
        \hrule height1pt
        \vspace{0.1cm}
        \hrule height2pt
        \vspace{0.3cm}
\end{minipage}

\begin{minipage}[c]{0.15\textwidth}
    \vspace{0.5cm}
    \hspace{0.5cm}
    \vrule width1pt height15cm 
    \hspace{0.05cm}
    \vrule width2pt height15cm
\end{minipage}
\hfill % Espacio horizontal entre las cajas
\begin{minipage}[c]{0.8\textwidth}
    \centering
        \textbf{\large Optimización dinámica de sistemas de reacción mediante cálculo fraccionario: Parte 1. Enfoque general para la formulación de modelos cinéticos fraccionarios de orden superior}\\[1cm]
        \text{\large Por}\\[1cm]
        \text{\Large Luis Raúl Barajas Villarruel}\\[1cm]
        \textit{\large Avance de tesis presentado al Departamento de Ingeniería Química como parte del proceso para obtener el grado de:}\\[1cm]
        \textbf{\Large Doctor en Ciencias}\\[1cm]
        \text{\large Dirigido por:}\\[1cm]
        \text{\Large PhD. Vicente Rico Ramírez}\\[1cm]
\end{minipage}

\begin{minipage}[c]{0.15\textwidth}
  \centering
  \includegraphics[width=\textwidth]{01_Figuras/00_Celaya.png}
\end{minipage}
\hfill % Espacio horizontal entre las cajas
\begin{minipage}[c]{0.8\textwidth}
    \centering
        \mdseries{\large Celaya, Guanajuato, \the\year}
\end{minipage}

\end{titlepage}
\restoregeometry
%%%%%%%%%%%%%%%%%%%%%%%%%%% RESUMEN %%%%%%%%%%%%%%%%%%%%%%%%%%%%%
\thispagestyle{empty}
\newgeometry{left=2.5cm,right=2.5cm,top=2.5cm,bottom=2.5cm}
\addcontentsline{toc}{chapter}{Resumen}
\begin{center}
\textsc{\Huge Resumen}\\[0.5cm]
\textbf{\large Optimización dinámica de sistemas de reacción mediante cálculo fraccionario: Parte 1. Enfoque general para la formulación de modelos cinéticos fraccionarios de orden superior}\\[0.25cm]
\text{\large Por: Luis Raúl Barajas Villarruel}\\[0.5cm]
\end{center}

lol

\begin{flushright}
    \textit{\large Dirigido por:}\\
    \textit{\large Vicente Rico Ramírez, PhD.}
\end{flushright}


\restoregeometry
%%%%%%%%%%%%%%%%%%%%%%%%%%% TABLA DE CONTENIDO %%%%%%%%%%%%%%%%%
\tableofcontents

\newpage
\pagenumbering{arabic}
\chapter{Introducción}

\citet{NewmanGirvan2004}

\newpage
\bibliographystyle{elsarticle-num-names}
\bibliography{03_Bibliografias.bib}
\end{document}


\end{document}